\section{Method}
\label{sec:method}

\subsection{Datasets and Common Processing Pipeline}

We evaluate on NF-UNSW-NB15 and NF-ToN-IoT, the NetFlow-standardised variants of their
respective intrusion-detection corpora \citep{moustafa2015unsw,sarhan2021netflow}. One
pipeline processes both, and every architectural choice is shared between them. Only two
quantities are permitted to differ, and both are selected on validation folds: the
uncertainty threshold $k$ and the feedback injection scale. This restriction is deliberate.
A cross-dataset comparison in which the two ladders use different components measures the
components, not the datasets.

The pipeline converts flow records into a graph, derives node and edge attributes, trains an
edge-level structural model, builds label-free semantic representations for the same edges,
and combines the two modalities through adaptive fusion. Every classification decision is
edge-level: a flow, or an aggregate of flows sharing a grouping key, is an edge, never a node.

\subsection{Directed Multigraph Construction}

Flows are aggregated into a directed multigraph keyed on

\[
(\text{source IP},\ \text{destination IP},\ \text{attack type}).
\]

Parallel edges therefore exist between an ordered IP pair whenever those flows carry
different attack types. Nodes are the distinct IP addresses. The resulting graphs are small:
NF-UNSW-NB15 yields 49 nodes and 656 edges, NF-ToN-IoT 1{,}501 nodes and 2{,}127 edges.

Each node carries ten centrality columns: betweenness, PageRank, degree, closeness,
eigenvector, $k$-core, $k$-truss, global betweenness, global PageRank, and modularity
vitality. These are computed on an unweighted IP-to-IP graph with self-loops removed.
Betweenness (normalised), PageRank, closeness and eigenvector centrality are computed on the
directed graph. Degree is \emph{in-degree} centrality, and $k$-core and $k$-truss are computed on
the undirected projection. Modularity vitality is the leave-one-node-out drop in modularity
under a greedy modularity partition of the undirected projection. Eigenvector centrality and
the truss decomposition fall back to constants when they fail to converge. A node's feature
vector is the mean of these values over all incident edges, taking the node in both its
source and destination roles.

One property of this feature set should be stated rather than left for a reader to discover.
On NF-UNSW-NB15 the ``global'' variants are populated from the same computation as their
local counterparts, so \texttt{global\_betweenness} and \texttt{global\_pagerank} are exact
duplicates of \texttt{betweenness} and \texttt{pagerank}. The ten columns span eight distinct
measures. On NF-ToN-IoT, whose centralities were computed upstream, the global variants
differ from the local ones. We report ten input channels because that is what the model
receives, but the effective rank of the node features is lower on NF-UNSW-NB15.

Each edge carries five attributes. Three are numeric: the number of constituent flows
(\texttt{flow\_count}), their summed byte count (\texttt{total\_bytes}), and their mean
duration (\texttt{avg\_duration}). Two are categorical: the modal protocol and the modal
port, with ties broken toward the smaller value.

Because attack type is part of the aggregation key, the label participates in defining edge
identity. This is a property of the protocol, not a bookkeeping detail. Two flows between the
same IP pair that differ only in their label become two distinct edges, and a deployed system
observing unlabelled traffic could not perform this separation. The scores reported in this
work therefore compare representations and mechanisms under a controlled protocol. They are
not deployment-valid detection estimates, and we do not present them as such.

\subsection{Edge Feature Encoding}

The encoding separates quantities from identifiers. Columns 0--2 receive a $\log(1+x)$
transform, then Z-scoring. Columns 3--4 are rewritten as contiguous vocabulary indices and
consumed by learned embedding tables of width 4 (protocol) and 32 (port). Concatenation gives
a $3 + 4 + 32 = 39$-dimensional encoded edge vector. The vocabularies are small: 7 protocols
and 169 ports on NF-UNSW-NB15, 4 and 235 on NF-ToN-IoT.

This corrects a specific defect in an earlier version of the pipeline, which passed protocol
and port through as raw magnitudes. Port then entered the network at roughly $11{,}000\times$
the scale of every normalised column, and the ordering it imposed is meaningless: port 443 is
not ``more'' than port 80. The embedding tables let the model represent protocol and port
identity without inheriting the arithmetic of their indices.

The Z-score statistics and the two vocabularies are fitted once over all edges, before the
cross-validation split. This is transductive, and we flag it explicitly. It is standard for a
single fixed graph, and it transfers no label information, but it is not the same as fitting
scalers inside each training fold.

\subsection{Structural Branch}

Each structural encoder is a two-layer GATv2 network \citep{brody2022gatv2}. The first layer
uses eight heads of width $64/8 = 8$, concatenated to 64. The second uses a single head of
width 64. Both layers receive the 39-dimensional encoded edge vector as an attention input and
add no self-loops. Each layer applies, in order, the GATv2 convolution, batch normalisation,
an ELU nonlinearity, and a residual connection. Dropout at $p = 0.2$ is applied before the
residual in the first layer only. The first layer's residual passes through a linear
projection from the 10 input channels to 64, and the second's is an identity. The final node
embedding is 64-dimensional.

We use GATv2 rather than GAT because of the static-attention limitation
\citep{brody2022gatv2}: in GAT the ranking a node induces over its neighbours can be
independent of the query node, so attention degenerates into a global neighbour scoring.
GATv2 applies its attention MLP after the concatenation of both endpoints, which restores
query-dependent ranking. On a graph with 49 nodes and heavy parallel-edge structure, that
distinction is not cosmetic.

Three encoders run in parallel: \emph{temporal}, \emph{content} and \emph{behavioural}.
They are architecturally identical and differ only in a fixed diagonal reweighting of the
encoded edge vector, given in Table~\ref{tab:focus}. The protocol weight is broadcast across
all 4 protocol-embedding dimensions and the port weight across all 32, so each variant keeps
the same relative emphasis it would have had on the five raw columns. The reweighted edge
vector is used both inside the attention computation and in the edge representation that
follows, so a variant's emphasis shapes its message passing as well as its readout.

\begin{table}[t]
\centering
\caption{Per-variant edge focus weights. Applied as a fixed diagonal scaling of the encoded
edge vector. The protocol and port weights are broadcast across their embedding blocks.}
\label{tab:focus}
\begin{tabular}{lccccc}
\toprule
Variant & flow\_count & total\_bytes & avg\_duration & protocol & port \\
\midrule
temporal    & 1.25 & 0.75 & 1.50 & 0.50 & 0.25 \\
content     & 0.50 & 1.50 & 0.50 & 1.00 & 1.25 \\
behavioural & 1.25 & 0.75 & 1.00 & 0.75 & 1.50 \\
\bottomrule
\end{tabular}
\end{table}

Within a variant, the representation of a directed edge $e = (u,v)$ concatenates the two
endpoint embeddings, their elementwise product, their elementwise absolute difference, and the
variant's reweighted edge vector $\tilde{\mathbf{x}}_{e}$:

\[
\mathbf{z}_{e}
=
\left[\,
\mathbf{h}_{u},\;
\mathbf{h}_{v},\;
\mathbf{h}_{u}\odot\mathbf{h}_{v},\;
\left|\mathbf{h}_{u}-\mathbf{h}_{v}\right|,\;
\tilde{\mathbf{x}}_{e}
\,\right]
\in \mathbb{R}^{4\cdot 64 + 39} = \mathbb{R}^{295}.
\]

The product and absolute-difference terms matter here. Endpoint concatenation alone cannot
separate parallel edges, and 58.7\% of NF-UNSW-NB15 edges share an IP pair with at least one
other edge. The elementwise terms are symmetric functions of the endpoints and do not fix
this on their own, but $\tilde{\mathbf{x}}_{e}$ does, by carrying the edge's own features into
the classifier.

Each variant projects $\mathbf{z}_{e}$ to 64 dimensions through a linear layer and ReLU. The
three projections are concatenated and passed through a linear layer to 64, ReLU, and dropout,
yielding the fused 64-dimensional edge embedding consumed by every downstream stage. A final
linear layer emits class logits. Each variant additionally carries its own classification head
, a two-layer MLP over $\mathbf{z}_{e}$ with a 64-unit hidden layer, ReLU and dropout,
whose loss is added to the objective with weight $0.30$. The auxiliary term keeps each variant
directly supervised rather than letting the fusion layer absorb all gradient.

\subsection{Semantic Branch}

The second modality describes each edge in controlled natural language. The continuous edge
quantities are first discretised into five rank-based quantile bands, computed globally across
all edges, and the band label rather than the raw number enters the sentence. A template then
renders endpoints, connection frequency, transferred volume, per-flow byte magnitude,
duration band, protocol name, and a coarse port role (web, remote-access, well-known,
registered, or ephemeral) into one sentence per edge, of the form: \emph{``Observed traffic
from source \dots to destination \dots with \dots connection frequency, transferring \dots as
\dots, \dots, over \dots, to a \dots port \dots''}

Sentences are encoded by CySecBERT \citep{bayer2022cysecbert}, a BERT-class encoder
domain-adapted to cybersecurity text. It is frozen throughout. No gradient reaches it. Inputs
are tokenised with padding and truncation at 128 tokens, and the final hidden states are
mean-pooled with padding positions masked out, giving one 768-dimensional vector per edge.
We use it strictly as a sentence encoder. It is not prompted, not asked to reason, and not
generative, and we avoid describing it as an LLM for that reason.

The embedding is consumed in two ways, and the distinction matters for everything that
follows. The \emph{embedding-only} route fits no classifier on it: a whitened prototype
scorer (described under the feedback stage below) assigns each edge to the nearest class
mean in a fold-local whitened space. This is the semantic rung of the ladder and the
semantic input to AGAF, and it measures what the encoder alone knows. The \emph{trained
head} route fits, within each fold, a small classifier on the training-fold embeddings: a
per-fold standard scaler, then a 768--128--128--$C$ multilayer perceptron with ReLU and
dropout $0.2$, trained with the same class-weighted focal loss, optimiser and early-stopping
rule as the structural branch. Each fold's head is applied to every edge and its logits are
stored: on that fold's training and validation edges the logits are in-fold, on its test
edges they are out-of-fold, so the pooled test score of the head is a legitimate out-of-fold
number. That head is the consultant the feedback loop uses, and, evaluated on its own, it is
also the strongest single-modality baseline in this study. We report it in every table for
that reason. A cross-fitted variant, in which training-edge logits come from an inner
cross-validation so that the loop trains on advice no more reliable than the advice it is
scored with, was implemented and measured. It scored lower on both datasets and is reported
in Section~\ref{sec:results}, and the in-fold head is retained.

Label freedom is enforced in code, not by convention. Before any sentence reaches the encoder,
an assertion scans every generated sentence for every attack term present in the dataset and
raises on the first match, naming the offending index and sentence. The pipeline halts. This
is what makes the semantic branch an independent view of the edge rather than a relabelling
of it.

\subsection{Adaptive Gated Attention Fusion}

AGAF projects each modality into a shared 128-dimensional space through a linear layer, batch
normalisation and ReLU, so 64 to 128 for the structural embedding and 768 to 128 for the semantic
one. A gate consumes the concatenated projections and emits two logits, softmaxed at
temperature 1.0 into weights $(\alpha_{g}, \alpha_{l})$ that convexly combine them:

\[
\mathbf{f}_{e} = \alpha_{g}(e)\,\mathbf{g}_{e} + \alpha_{l}(e)\,\mathbf{l}_{e},
\qquad
\alpha_{g}(e) + \alpha_{l}(e) = 1 .
\]

The weights are per-edge, so the balance between structure and semantics is free to vary
across the graph rather than being fixed at concatenation time. A two-layer MLP over
$\mathbf{f}_{e}$ produces class logits. Both inputs are standardised before entering the
module. The scalers are held in the trainer and persisted separately, so the module's
\texttt{state\_dict} stays free of fitted statistics.

One constraint on this stage is load-bearing. AGAF consumes \emph{per-fold out-of-fold}
structural embeddings only. It never receives embeddings from a GNN fitted on all edges,
because such a model has already seen the test edges' labels, and those labels would reach
the fusion classifier through its own input features. That is a leak no split downstream of it
could detect. For each held-out fold, the structural embeddings supplied to AGAF come from a
model that never saw that fold's labels. Earlier runs in this project violated this rule, and
the reproduction entry point now enforces it.

\subsection{Uncertainty-Guided Semantic Feedback}

The feedback stage routes semantic information to the edges the structural model is least sure
about, and only there.

\paragraph{Selection.} For each edge we compute the Shannon entropy, in nats, of the GNN's
predictive distribution, and flag the edges above the $(1 - k/100)$ quantile. The threshold is
calibrated once and reused, so the flagged set is a fixed-size top-$k\%$ rather than a moving
target. We select $k$ on validation folds only, over $[15, 35]$ at three training seeds:
$k = 29.0$ on NF-UNSW-NB15 and $k = 16.0$ on NF-ToN-IoT. The validation curves are flat, with
every candidate inside the noise band of the mean, so these values are the argmax of a flat
curve rather than a tuned optimum, and we say so wherever they appear.

\paragraph{Consultation.} The loop consults the per-fold trained head described in the
semantic branch. For every flagged edge it reads that fold's stored head logits, a
$C$-dimensional class judgement computed from the edge's sentence alone. The judgement does
not depend on the graph model's current belief about the edge, so the coupling is
bidirectional in routing but not in content.

The embedding-only alternative, used by the semantic rung and by AGAF, is a whitened
prototype scorer. Within each fold we centre the training-fold embeddings, fit a ZCA whitener
by eigendecomposition of their covariance with a floor on small eigenvalues (necessary
because the training fold has fewer rows than the 768 embedding dimensions), and take the
class-wise mean in the whitened space as that class's prototype. An edge is whitened with the
same fold-local transform and scored by cosine similarity against the prototypes. Whitening
is not incidental: an earlier L2-normalise-average-renormalise scheme collapsed per-class
spread until random and real advice became indistinguishable.

The rule that keeps the comparison fair is stated once here and applied everywhere: every
rung uses the same encoder, the same graph, the same folds and the same training seeds. The
loop additionally trains a classification head on the semantic embeddings, and that head is
reported alone. The loop was first built with the prototype scorer as its consultant, on the
argument that a consultant the loop already contains would only be echoed. That
configuration, and five attempts to make its advice usable, are reported as negative results
in Section~\ref{sec:results}. The trained head replaced it after those measurements.

\paragraph{Gating.} Being uncertain in the structural view does not make an edge a good
candidate for semantic advice. Among the flagged edges we keep only the most confident half by
the consultant's own maximum class probability, a semantic confidence fraction of $0.50$. The
effective feedback rate is thus $14.5\%$ of edges on NF-UNSW-NB15 and $8.0\%$ on NF-ToN-IoT.

\paragraph{Injection.} Accepted advice is injected into the \emph{edge representation}
consumed by the next graph update. The consultant's class logits are projected linearly to the
39-dimensional edge width, scaled by a learned strength, written into a full-graph bias tensor
that is exactly zero on unflagged edges, and added to the reweighted edge vector:

\[
\tilde{\mathbf{x}}_{e} \leftarrow \tilde{\mathbf{x}}_{e} + s \cdot \mathbf{b}_{e},
\qquad
\mathbf{b}_{e} = \mathbf{0} \ \text{for unflagged } e ,
\]

with injection scale $s = 2.0$ on NF-UNSW-NB15 and $s = 20.0$ on NF-ToN-IoT, both selected on
validation folds from $\{0.5, 1, 2, 5, 10, 20\}$ at three seeds. As with $k$, the curves are
flat, and on NF-ToN-IoT the selected value sits on the upper boundary of the swept range,
which we record rather than hide. Because $\tilde{\mathbf{x}}_{e}$ enters both attention and readout, the
injected signal propagates structurally rather than only shifting a final logit. The scale is
re-selected whenever the edge encoding or the injection path changes, since it is expressed in
the units of that encoding.

\paragraph{Iteration.} Predictions are recomputed, entropies re-derived, and the cycle repeats.
We measure \emph{churn}, the fraction of edges whose predicted class changed since the previous
iteration, and stop when churn falls below $0.01$ or after three iterations, whichever comes
first.

\paragraph{Output fusion.} The loop's prediction is not the graph head's logits after the
last iteration. The refined edge embedding and the consultant's logits are each projected to
128 dimensions and mixed by a per-edge gate $g_e \in (0,1)$ that reads the entropy and
maximum probability of both the graph head and the consultant, so it can route towards
whichever view is confident on that edge. A two-layer classifier over the gated concatenation
produces a correction, and the final logits are that correction plus the gated residual
$(1-g_e)\,\mathbf{z}^{\text{gnn}}_e + g_e\,\mathbf{z}^{\text{head}}_e$. The consultant's
own judgement is therefore a floor the loop can fall back to on edges where the graph view is
unsure. This stage has the same projection width as AGAF, so the loop and AGAF differ in
where semantics enter (inside the graph update, or only at the output) and not in fusion
capacity. Switching the stage off, together with the injection, yields the \texttt{head\_only}
control used throughout Section~\ref{sec:results}: the same network, trained in the same
code path, with no semantic input at all.

\paragraph{An alternative that was implemented and measured.} We also built a variant that adds
the projected semantic signal to GATv2's unnormalised attention logits, before the per-edge
softmax, injecting into the attention mechanism rather than the edge representation. It is
zero-initialised, so an untrained bias reproduces stock GATv2 exactly. It is evaluated under
the identical protocol and reported in Section~\ref{sec:results} as a negative result. We
retain it in the paper because the two variants isolate \emph{where} in the network a semantic
signal must land to have any effect at all, which is a question the positive result
alone does not answer.

\subsection{Baseline Protocol}

We compare against re-implementations of two published graph-based intrusion detectors,
E-GraphSAGE \citep{lo2022egraphsage} and TE-G-SAGE \citep{tegsage2025}, retrained on the same
aggregated edges, the same labels and the same folds. Each baseline keeps its own published
featurisation rather than ours. E-GraphSAGE standard-scales all five NetFlow columns,
TE-G-SAGE applies log1p, standard scaling, correlation pruning at 0.995, and one-hot
categoricals with rare-category folding.

Each baseline is run in three configurations. \texttt{as\_published} reproduces the released
implementation. \texttt{refit} re-selects capacity and schedule on our validation folds only.
\texttt{plus\_node\_features} is deliberately non-faithful: it hands the baseline our ten
centrality node features, isolating how much of any gap is representation rather than
architecture. We label it as non-faithful wherever it appears. It is not E-GraphSAGE and it is
not TE-G-SAGE.

Five deviations from the published specifications were unavoidable at this graph scale, and we
record them rather than absorb them. (D1) TE-G-SAGE's chronological 60/30/10 split cannot be
reconstructed, because aggregation collapses per-flow timestamps. Our stratified folds are used
for all models. (D2) TE-G-SAGE's neighbour fanout of $(25, 15)$ and batch size of 4096 both
exceed the entire graph, so training is full-neighbourhood and full-batch. (D3) Both released
implementations train for a fixed epoch count with no validation split and no early stopping
, 4999 epochs for E-GraphSAGE and 20 for TE-G-SAGE, which is incompatible with a protocol
that selects on validation folds. Baselines are given 300 epochs with patience 25 on validation
macro-F1 and the best state restored. (D4) TE-G-SAGE's \texttt{rare\_min\_freq} of 50 removes
almost every port category at our scale. It is applied as published in \texttt{as\_published}
and swept in \texttt{refit}. (D5) E-GraphSAGE's published Eq.~4 aggregates edge features alone,
but the authors' released code builds the message as $W_{\text{msg}}([\mathbf{h}_u \Vert
\mathbf{e}_{uv}])$. We follow the code, since the published numbers came from it.

One design variable is internal to the loop rather than a rung of the ladder: the identity of
the semantic consultant. The loop was first built around the whitened prototype scorer, which
is also what the semantic rung uses, and later around a per-fold classification head trained on
the same frozen embeddings. Both versions are reported, at the same three seeds, at the same
fold partition, and with the injection scale unchanged. Only the selected entropy percentile
differs between them, at 31 against 29 on NF-UNSW-NB15 and 25 against 16 on NF-ToN-IoT.

Two asymmetries remain and should be weighed when reading the comparison. Our rungs train
under a focal objective with patience 40, whereas the baselines use class-weighted
cross-entropy with patience 25, each keeping its own published loss. And E-GraphSAGE
classifies an edge from $[\mathbf{h}_u \Vert \mathbf{h}_v]$ alone, so parallel edges between
one IP pair are mathematically indistinguishable to it. That covers 58.7\% of NF-UNSW-NB15 edges and
7.9\% of NF-ToN-IoT edges fall in such groups. That is an interaction between its architecture
and our aggregated representation, not evidence that E-GraphSAGE is a weak detector, and we do
not read it as one.

\subsection{Evaluation and Training Protocol}

We use stratified five-fold cross-validation over edges with partition seed 42, represented as
boolean masks. Within each fold the held-out fifth is the test partition, and the remaining
four-fifths are split 75/25 into training and validation with seed $42 + \text{fold index}$,
giving approximately $60/20/20$. Test masks are used once, for prediction.

The primary metric is pooled out-of-fold macro-F1: each edge is predicted only while held out,
and the five sets of out-of-fold predictions are pooled before a single macro-F1 is computed.
Model selection uses validation folds exclusively. The uncertainty threshold $k$ and the
injection scale were selected per dataset by a validation sweep ($k \in [15, 35]$; scale
$\in \{0.5, 1, 2, 5, 10, 20\}$) averaged across three seeds (42, 1, 2). The remaining
architectural hyperparameters were fixed a priori and not swept. We note this as a limitation:
$k$ was selected on rotating validation folds rather than by nested cross-validation, so the
final estimate is not fully unbiased with respect to that one choice.

Every rung is trained at three training seeds (42, 1, 2) with the fold partition held fixed at
the seed-42 split, so all seeds score the same held-out edges and differ only in
initialisation, dropout and data order. We report the mean and standard deviation of pooled
out-of-fold macro-F1 over the three seeds. The standard deviation therefore measures
training-seed variance at a fixed partition, not total uncertainty, since fold-partition
variance is not measured. Pairwise differences between rungs use a two-level bootstrap with
2{,}000 iterations: each iteration draws a training seed independently for each side and
resamples edges with replacement, so the interval reflects both sources of variation. A
seed-matched bootstrap (seed 42 against 42, 1 against 1, 2 against 2) is reported as a
secondary check. The decision rule was fixed before the runs: a difference is called
\emph{separated} only if its two-level 95\% interval excludes zero, and \emph{not separated}
otherwise; we do not describe point-estimate orderings as significant.  % numcheck: allow -- the sentence refuses the word
The semantic rung is
deterministic given the folds and has zero seed variance, so it enters comparisons as a fixed
side. Comparisons against re-trained published baselines use the same procedure with three
baseline seeds on the other side.

Training uses AdamW with learning rate $10^{-3}$ and weight decay $5\times10^{-4}$, at most 300
epochs, early stopping with patience 40 on validation macro-F1 with the best state restored,
and gradient clipping at norm 1.0. The objective is focal loss with $\gamma = 2.0$ and $\alpha$
set to inverse-frequency class weights computed on the training fold, normalised to sum to the
class count, and classes absent from a training fold receive weight zero. Focal loss down-weights
already-correct examples, which matters when rare attack classes hold 12--35 edges against a
benign majority. We do not additionally oversample: combining oversampling with these weights
double-corrects the imbalance and collapses the fusion stage.

Every run exports \texttt{OMP\_NUM\_THREADS=1}. Graph-attention message passing sums neighbour
contributions in thread-completion order, and floating-point addition is not associative, so
thread count alone moves pooled macro-F1 by roughly 0.012 at a fixed seed, larger than
several of the effects this paper reports. Semantic encoding is likewise pinned to a single
device, because MPS and CPU produce slightly different embeddings and every downstream number
inherits the difference.

Two NF-ToN-IoT classes are smaller than the five-fold minimum: \texttt{dos} with 4 edges and
\texttt{ransomware} with 3. They remain in the graph and continue to participate in message
passing, but are excluded from the reported metric, so NF-ToN-IoT macro-F1 is computed over 8
classes and NF-UNSW-NB15 over 10. The two are therefore not comparable in absolute level, and
we compare only the shape of the ladder across datasets, never the values.
